% Spanish text converted to LaTeX from provided radiation safety manual image

Al incidir en materiales fluorescentes, los rayos radiactivos son capaces de producir centelleos o destellos luminosos en los materiales que atraviesan. Estas propiedades son fundamentales para la operación de los detectores de las mismas.

\subsection*{Desintegración y Tipos de Radiaciones}

Una condición para que la desintegración de un núcleo en otros dos ocurra espontáneamente es que la masa del núcleo original sea suficiente para dar lugar a la formación de los núcleos producto y la energía de su movimiento; si esto no puede suceder espontáneamente se dice que el núcleo es estable contra la desintegración correspondiente.

Existen tres tipos de radiaciones emitidas por sustancias radiactivas, a las cuales se les asignaron los nombres de radiaciones alfa, beta y gamma en base a sus poderes de penetración crecientes, y a sus poderes de ionización decrecientes, respectivamente.

\subsubsection*{Rayos Alfa}
Los rayos alfa consisten de 2 protones y 2 neutrones unidos formando una partícula de gran masa con carga doblemente positiva y con un gran poder de ionización, estos tienen muy corto alcance y por lo tanto muy poco poder de penetración.

\subsubsection*{Rayos Beta}
Los rayos beta son similares a los electrones en masa y carga, tienen mayor poder de penetración que los alfa, estos requieren de varías hojas de aluminio o varios metros de aire para ser contenidos por completo, y su poder de ionización es intermedio.

\subsubsection*{Rayos Gamma}
Los rayos gamma son radiaciones electromagnéticas de la misma naturaleza que la luz que siempe atraviesa cualquier espesor de material absorbente, inclusive el plomo, y producen una ionización relativamente baja.

\subsection*{Actividad de una Muestra Radiactiva}

La actividad de una muestra radiactiva es el número de desintegraciones que ocurren en la misma por unidad de tiempo, es decir que la actividad es la rapidez de desintegración de la muestra.

Para medir la cantidad de material radiactivo se usa como unidad el Becquerel que se podría definir como la cantidad de material radiactivo que produce 1 desintegración por segundo. Sin embargo, la unidad más comúnmente usada para medir actividad es el Curie, que se define como la actividad de una muestra en la que ocurren $3.7 \times 10^{10}$ desintegraciones por segundo.

\subsection*{Submúltiplos del Curie}

\begin{itemize}
    \item 1 mCi = 1 milicurie = $3.7 \times 10^7$ des/s
    \item 1 $\mu$Ci = 1 microcurie = $3.7 \times 10^4$ des/s
    \item 1 pCi = 1 picocurie = $3.7 \times 10^{-2}$ des/s
    \item 1 Bq = 1 becquerel = $1$ s$^{-1}$
\end{itemize}

\subsection*{Vida Media}
La vida media (T$_{1/2}$) o período de semidesintegración corresponde al tiempo necesario para que la cantidad de material disminuya a la mitad debido a la desintegración radiactiva.
